\subsection{可三角化的自同态}

在本节中，我们关注如下情形：u 的最小多项式 \(p(\mathbf{X})\) 分裂为一次因子的乘积（于 K{[}X{]} 中），换言之（VII，第 33 页，推论 2），即 u 的所有特征值都属于 K 的情形。当 K 是代数闭域时，这一点尤其成立。VII，第 31 页的命题 3 立即给出：

命题7. --- 设 E 是交换域 K 上的有限维向量空间，且 u 是 E 的一个自同态，其特征值均属于 K。对 u 的每个特征值 \(\alpha\)，设 \(M_{\alpha}\) 为 E 的由那些存在整数 \(k \geqslant 1\) 使得 \((u - \alpha)^{k} \cdot x = 0\) 的元素 x 组成的向量子空间。则 \(M_{\alpha}\) 在 u 下封闭，向量空间 E 是诸 \(M_{\alpha}\) 的和，并且存在多项式 \(s_{\alpha} \in K[X]\)，使得对所有 \(x \in E\)，x 在 \(M_{\alpha}\) 中的分量等于 \(s_{\alpha}(u) \cdot x\)。

子模 \(M_{\alpha}\) 作为 K{[}X{]}-模是有限生成的，故其零化子具有形式 \((\mathbf{X}-\alpha)^{r}\)；换言之，存在整数 \(r \geqslant 1\) 使得

\[
(u - \alpha) ^ {r} \cdot x = 0
\]

对所有 \(x \in \mathbf{M}_{\alpha}\)；\(u - \alpha\) 在 \(\mathbf{M}_{\alpha}\) 上的限制是一个幂零自同态。

仍假设 u 的特征值属于 K，我们现在把 VII，第 34 页的命题 6 应用于 u。多项式 \(p_{k}\) 无非就是 X - \(\alpha\)（当 \(\alpha\) 遍历 u 的特征值集合时），我们看到 E 是诸子空间 \(E_{i}\) 的直和，这些子空间在 u 下封闭、关于 u 循环，并且使得 \(u\) 在 \(E_i\) 上限制的最小多项式具有形式 \((X - \alpha)^m\)。设 \(E_i'\) 为一个 K{[}X{]}-模，与 \(E_i\) 相伴；则 \(E_i'\) 同构于模 K{[}X{]}/\((X - \alpha)^{m}\) 之一。现在，模 (X - \(\alpha\) ) \(^m\) 的剩余类中，来自元素 (X - \(\alpha\) ) \(^k\) ( \(0 \leqslant k \leqslant m - 1\) ) 的那些构成 K{[}X{]}/((X - \(\alpha\) ) \(^m\) ) 的一组 K-基（IV，第 11 页，推论），且

\[
\mathbf {X} (\mathbf {X} - \alpha) ^ {k} = \alpha (\mathbf {X} - \alpha) ^ {k} + (\mathbf {X} - \alpha) ^ {k + 1}
\]

给定 \(0 \leqslant k \leqslant m - 1\)；由此可得 \(\mathbf{E}_i\) 的维数为 \(m\)，且若 \(\alpha\) 是限制 \(u_i\)（即 \(u\) 在 \(\mathbf{E}_i\) 上的限制）的唯一特征值，则存在 \(\mathbf{E}_i\) 的一组基，使得 \(u_i\) 关于该基的矩阵是 \(m \times m\) 矩阵

\[
U _ {m, \alpha} = \left( \begin{array}{c c c c c c} \alpha & 0 & 0 & ... & 0 & 0 \\ 1 & \alpha & 0 & ... & 0 & 0 \\ 0 & 1 & \alpha & ... & 0 & 0 \\ \cdots & \cdots & \cdots & \cdots & \cdots & \cdots \\ 0 & 0 & 0 & ... & \alpha & 0 \\ 0 & 0 & 0 & ... & 1 & \alpha \end{array} \right)\tag{4}
\]

定义3. --- 对每个域 K、每个整数 \(m \geqslant 1\) 以及每个 \(\alpha \in K\)，矩阵 \(U_{m,\alpha}\) 称为阶为 m、特征值为 \(\alpha\) 的若尔当矩阵。

命题8. --- 设 E 是交换域 K 上的有限维向量空间，且 u 是 E 的一个自同态。则以下条件等价：

\begin{enumerate}
\def\labelenumi{(\roman{enumi})}
\item
  \(u\) 的特征值（在 \(\mathbf{K}\) 的某个代数闭扩张中）属于 \(\mathbf{K}\)；
\item
  存在 \(\mathbf{E}\) 的一组基，使得 \(u\) 关于该基的矩阵是下（相应地，上）三角的；
\item
  存在E的一组基，使得u在这组基下的矩阵是由若尔当矩阵构成的对角分块矩阵。
\end{enumerate}

由命题 7 及上述注记，我们有 (i) \(\Rightarrow\) (iii)；而断言 (iii) \(\Rightarrow\) (ii) 和 (ii) \(\Rightarrow\) (i) 是平凡的。

定义4. --- 满足命题 8 的条件 (i)、(ii) 和 (iii) 的自同态称为可三角化的。

特别地，若 K 是代数闭的，则 K-向量空间的每个自同态都是可三角化的。

对于矩阵，命题 8 蕴含：

推论。--- 设 U 是交换域 K 上的方阵，且 U 的所有特征值都属于 K；则存在一个与 U 相似的矩阵，它是若尔当矩阵的一个对角分块。

注记。--- 1) 由 VII，第 34 页，命题 6 可知，若 U 相似于由若尔当矩阵构成的一个对角分块矩阵 \((J_{k})\)，则 \(J_{k}\) 中形如 \(U_{m,\alpha}\)（对给定的 m 与 \(\alpha\)）者的个数由 U 唯一确定。

\begin{enumerate}
\def\labelenumi{\arabic{enumi})}
\setcounter{enumi}{1}
\tightlist
\item
  更一般地，如果 \(U\) 相似于若尔当矩阵的一个对角分块 \(U_{m_i,\alpha_i}\)，则 \(U\) 的相似不变量可以仿照 VII，第 25 页，注记 3 中给出的方法直接计算：把所有 \((X - \alpha_i)^{m_i}\)（\(\alpha\) 相同者）写在同一行，按指数的降序排列，并用 1 补齐，使每行的长度等于 \(U\) 的阶；这样，\(U\) 的相似不变量便按指标的降序、通过将同一列中的各项相乘而得到。例如，对于矩阵
\end{enumerate}

\[
\left( \begin{array}{c c c} 2 & 0 & 0 \\ 0 & 3 & 0 \\ 0 & 1 & 3 \end{array} \right)
\]

我们写出

\[
\begin{array}{l} (\mathbf {X} - 2), 1, 1 \\ (\mathbf {X} - 3) ^ {2}, 1, 1 \end{array}
\]

，相似不变量为 1、1 和 \((\mathbf{X} - 2)(\mathbf{X} - 3)^2\)。

注意到若尔当矩阵 \(U_{m,\alpha}\) 的最小多项式是 \((X-\alpha)^{m}\)，并且它等于其特征多项式，我们得到以下结果：

命题9. --- 如果方阵 U 相似于由若尔当矩阵构成的对角分块矩阵 \((U_{m_{i},\alpha_{i}})\)，那么 U 的最小多项式是这些 \((\mathbf{X}-\boldsymbol{\alpha}_{i})^{m_{i}}\) 的最小公倍数，而特征多项式是这些 \((\mathbf{X}-\boldsymbol{\alpha}_{i})^{m_{i}}\) 的乘积。

推论。--- 在命题 7 的记号下，子空间 \(M_{\alpha}\) 的维数是特征值 \(\alpha\) 作为 u 的特征多项式的根的重数。

